\section{Introduction}\label{sec:introduction}
Let $(g_{ij})_{1\le i<j\le n}$ be independent standard Gaussian random variables. For a fixed inverse temperature $\beta>1$, define
\[
 H_n(\sigma)=\frac{\beta}{\sqrt n}\sum_{1\le i<j\le n}g_{ij}\sigma_i\sigma_j,
 \qquad Z_n=2^{-n}\sum_{\sigma\in\{-1,1\}^n}e^{H_n(\sigma)},
 \qquad F_n=\log Z_n.
\]
The normalization of the spin sum changes $F_n$ by a deterministic constant and therefore has no effect on centered fluctuations.

\begin{theorem}\label{thm:main}
For every fixed $1<\beta<\infty$,
\[
 \sd(F_n)=n^{1/6+o(1)},\qquad \Var(F_n)=n^{1/3+o(1)}.
\]
More explicitly, for every $\eta>0$ there is $n_0(\beta,\eta)$ such that
\[
 n^{1/6-\eta}\le\sd(F_n)\le n^{1/6+\eta}
 \qquad(n\ge n_0(\beta,\eta)).
\]
Moreover, for every $\eta>0$,
\[
 \mathbb P\!\left(n^{1/6-\eta}\le|F_n-\mathbb EF_n|
                         \le n^{1/6+\eta}\right)\longrightarrow1.
\]
\end{theorem}
For the free-energy density $F_n/n$, the standard-deviation and typical-magnitude exponents are $-5/6$, and the variance exponent is $-5/3$. Multiplication by the fixed physical factor $-1/\beta$ has no effect on these exponents. Theorem~\ref{thm:main} concerns power exponents; it gives neither a limiting law nor a nonzero constant prefactor. The temperature is fixed throughout, with no uniform conclusion as $\beta$ approaches $1$ or infinity.

\subsection{History and significance}
The Sherrington--Kirkpatrick model \cite{SK1975} separates two questions:
the limiting free energy and its fluctuations between disorder samples.
Parisi introduced the hierarchy of order parameters underlying the
variational prediction for the first \cite{Parisi1979}.
Guerra and Toninelli \cite{GuerraToninelli2002} established the
thermodynamic limit by Gaussian interpolation, and Guerra's variational
bound \cite{Guerra2003} and Talagrand's matching result
\cite{Talagrand2006} identified its Parisi value. Panchenko subsequently
established the formula for general mixed $p$-spin models, including odd
interactions \cite{Panchenko2014}. These results determine the leading
extensive term; the fluctuations studied here occur on a smaller scale.

The prediction of that scale has a separate history. Kondor's replica
expansion near the critical temperature \cite{Kondor1983} is the starting
point of the calculation described by Parisi and Rizzo
\cite{ParisiRizzo2008,ParisiRizzo2009}. Crisanti, Paladin, Sommers, and
Vulpiani \cite{CPSV1992} predicted a scale $n^{-5/6}$ for the physical
free energy per spin in the low-temperature, zero-field phase. At fixed
$\beta$, this corresponds to $n^{1/6}$ for the log partition function.
Parisi and Rizzo developed the associated rare-deviation calculation
through the low-temperature phase. Inferring typical fluctuations from
that calculation requires an additional matching assumption, discussed
in \cite[Section~VI]{ParisiRizzo2009}.
Theorem~\ref{thm:main} establishes the predicted one-sixth exponent with
arbitrarily small power slack, for every fixed finite $\beta>1$.

At fixed high temperature $0<\beta<1$, order-one Gaussian fluctuations
go back to Aizenman, Lebowitz, and Ruelle
\cite{AizenmanLebowitzRuelle1987}. Chatterjee's work on disorder chaos
and superconcentration \cite[Theorem~1.5]{Chatterjee2009} gives
$\operatorname{Var}(F_n)=O_\beta(n/\log n)$, and his later lower-bound
method proves constant-scale nonconcentration
\cite[Theorem~2.4]{Chatterjee2019LowerBounds}.
Chen and Lam obtained a critical-temperature variance upper bound of
order $(\log n)^2$ \cite{ChenLam2019Critical}. At the critical point,
Du and Huang subsequently proved variance $\frac16\log n+O(1)$ and
a Gaussian limit \cite[Theorem~1.3]{DuHuang2026Critical}.
Critical-window limit theorems concern temperatures approaching the
transition
\cite{ChengLiuShaoXu2026CriticalWindow}.

For the fixed low-temperature regime considered here, Aronow and Lopatto
\cite[Theorem~1.4]{AL2026} prove, for all sufficiently large $n$,
\[
 c_\beta n^{4/15}\le \operatorname{Var}(F_n)
                         \le C_\beta n^{7/15}.
\]
Their positive-transform and deviation estimates are a direct
antecedent of the present argument.

The structure of the scalar Parisi minimizer is also essential.
Auffinger and Chen established origin-support and regularity properties
\cite{AuffingerChenProperties2015}, and proved uniqueness through a
stochastic-control representation \cite{AuffingerChenUniqueness2015}.
Jagannath and Tobasco developed variational support and self-consistency
conditions \cite{JagannathTobasco2017}. After Zhou's interval-support
result near the critical temperature \cite{Zhou2025}, Lopatto
\cite[Theorem~1.1]{Lopatto2026} established interval support $[0,q_\beta]$
and the corresponding density structure at every fixed $\beta>1$.
We use this theorem together with the support self-consistency in
\cite[Lemma~2.2]{Lopatto2026}. The strictly positive initial density
needed here is then derived from scalar
diffusion identities and a strict fourth-derivative estimate in
Section~\ref{sec:profile}.

\subsection{Proof strategy and method ancestry}
Set $T=\beta^2$. We work with the completed covariance
\[
 \mathbb E[H_n^c(\sigma)H_n^c(\tau)]=\frac{nT}{2}R(\sigma,\tau)^2,
 \qquad R(\sigma,\tau)=\frac{\sigma\cdot\tau}{n}.
\]
Here $H_n^c=H_n+D$, where $D$ is an independent common Gaussian of
variance $T/2$, and $F_n^c=F_n+D$. This completion gives the displayed
quadratic covariance; the added noise is removed in the final step.

The quantitative target is a positive Laplace window for
$\log\mathbb E\exp\{s(F_n^c-\mathbb EF_n^c)\}$ when $s$ exceeds
$n^{-1/6}$ by an arbitrarily small fixed power of $n$
(Proposition~\ref{fl:laplace-window}). To compare these transforms,
we reveal Gaussian covariance in stages and integrate a stage at mass
$u$ by the operation $X\mapsto u^{-1}\log\mathbb E e^{uX}$, with
expectation at $u=0$. Nondecreasing paths specify the cumulative
interaction and external-field covariance. Copies of the spin system
can share selected stages of this revelation; their prescribed common
prefixes form the finite trees used in the proof.

This hierarchical framework belongs to the lineage of Parisi recursion,
Ruelle's cascades \cite{Ruelle1987}, and Gaussian interpolation
\cite{GuerraToninelli2002,Guerra2003}. The cavity variational principle
of Aizenman, Sims, and Starr \cite{AizenmanSimsStarr2003} is a further
antecedent for comparisons that remove or add sites. We construct the
finite positive transition laws used here directly.

The closest analytic predecessor is Aronow and Lopatto's combination of
a pinned Brascamp--Lieb variance estimate, optimized covariance
interpolation, and positive moments
\cite[Propositions~2.2--2.3 and Section~3]{AL2026}.
The extensions proved here concern paired interaction and field clocks,
local comparisons, fixed higher moments, and grouped replica tests at
the fluctuation scale. The continuous inverse-Hessian inequality is due
to Brascamp and Lieb \cite{BL1976}. Its application after pinning one
terminal spin configuration is proved in Section~\ref{sec:foundations};
the resulting moment estimate is for the full unconditioned planted
path, not for an arbitrary fixed prefix.

A global replacement of both hierarchical pressures by their scalar
Parisi limits loses too much at this scale. Instead,
Section~\ref{sec:comparisons} optimizes short covariance comparisons
over a nondecreasing overlap quantile. The resulting error consists of
overlap variance and the mismatch between that quantile and the mean
overlap. Under comparability of mass and overlap above the scale being
tested, projection and pinned moment estimates make these errors small
in every fixed moment at selected split levels. The polynomial saving
can be chosen independently of that fixed moment order.

Section~\ref{sec:cavity} turns these estimates into a scalar description
of the mean overlap. Removing one spin produces a susceptibility
correction that is too large under an absolute term-by-term bound.
Additional test replicas give a finite system of identities whose
iteration cancels that correction. For a block already containing
observed replicas, the signed attachment contributions are summed
before absolute values are taken. The number of test factors and the
iteration depth remain fixed as $n$ grows. This gives the cubic
one-site closure of Proposition~\ref{prop:cavity-closure}, including
the uniform version needed for moving split locations.

Section~\ref{sec:profile} uses that closure and the comparison's first
variation to determine its small-mass profile. A limiting scalar
equation and a nonnegative variational potential exclude gaps in the
rescaled overlap distribution, forcing a positive linear profile.
The estimates hold in both mass and overlap coordinates, so an interval
of negligible mass cannot conceal a gap of appreciable overlap length.
The same section identifies macroscopic limits with the unique Parisi
minimizer; its control proof follows the strict-convexity method of
Auffinger and Chen \cite{AuffingerChenUniqueness2015}.

The local arguments assume regularity above the scale being tested.
Section~\ref{sec:regularity} removes this assumption by finite
propagation. A failure is examined at a scale with a regular outer
region. Sufficient comparison strength there already contradicts the
direct profile result. In the remaining case, a stronger local
comparison is constructed. If its error bound holds, agreement with
the original comparison yields the same contradiction; if it fails,
the next comparison has an error budget larger by a fixed power of $n$.
A fixed number of repetitions exhausts the admissible range.
Figure~\ref{fig:proof-order} records this alternative.

\begin{figure}[htbp]
\centering
\begin{tikzpicture}[
 box/.style={draw,rounded corners=2pt,align=center,font=\small,
             text width=5.55cm,inner sep=5pt,minimum height=1.05cm},
 branch/.style={box,minimum height=3.3cm},
 flow/.style={-{Stealth[length=2mm]},thick}]
\node[box] (bad) at (0,0)
 {Selected bad parent scale\\Regular outer region};
\node[box] (child) at (7.1,0)
 {Admissible child\\Stronger local comparison};
\node[branch] (success) at (0,-3.0)
 {Child error bound holds\\[2pt]
  Agreement with parent\\[4pt]
  Small scale: outer regularity,\\
  then the direct profile\\[2pt]
  Macroscopic scale: Parisi profile\\[4pt]
  Contradiction at the parent scale};
\node[branch] (failure) at (7.1,-3.0)
 {Child error bound fails\\[4pt]
  Select a strength and minimizer\\
  with a bad child profile\\[4pt]
  Larger error-budget parameter};
\draw[flow] (bad.east) -- (child.west);
\draw[flow] (child.south west) -- (success.north east);
\draw[flow] (child.south) -- (failure.north);
\draw[flow] (failure.east) -- (10.65,-3.0) -- (10.65,1.2)
 -- node[above,font=\small] {Repeat at the selected child strength}
 (0,1.2) -- (bad.north);
\end{tikzpicture}

\caption{A selected bad parent scale has a regular outer region. The diagram
concerns the remaining weak-strength case; a direct profile test
already excludes sufficient local strength. Child-error success gives
agreement. At a scale tending to zero, this supplies child outer
regularity for the small-scale direct test; at a macroscopic scale,
the Parisi profile applies instead. Either profile contradicts the
parent failure (Lemma~\ref{reg:successful-child}). Child-error failure
selects a strength and minimizer, and then the next bad scale. Each
child raises the error-budget parameter by a fixed power of system
size. The covariance budget preserves admissibility, while this growth
exhausts the allowed error range after a fixed number of steps
(Proposition~\ref{reg:finite-propagation}). This propagation is separate
from the local finite-degree cavity recursion.}
\label{fig:proof-order}
\end{figure}

Finally, Section~\ref{sec:fluctuations} applies the resulting comparisons
to obtain the positive Laplace window. A positive transform alone does
not control variance or typical magnitude. Following the exponential
augmentation argument of \cite[Proposition~5.1]{AL2026}, we add an
independent rate-one exponential variable to $F_n^c$. Gauge symmetry
and Pr\'ekopa's marginal theorem \cite{Prekopa1973} give a log-concave
density. Its moment and small-ball bounds convert the Laplace window
into the standard-deviation and typical-magnitude assertions. Removing
the added exponential and Gaussian variables completes the proof.

\subsection{Conventions}
All constants may depend on the fixed $T$. Small auxiliary exponents are chosen first and held fixed as $n\to\infty$. Assertions along test sequences permit passage to subsequences, but never an $n$-dependent iteration depth. We write $A\lesssim B$ for an inequality with a constant independent of $n$, and $A\asymp B$ when both inequalities hold. Dependence on a fixed moment order is indicated by a subscript. A polynomial saving means a factor $n^{-c}$ for some fixed $c>0$; logarithmic losses are absorbed only after this saving has been established.
